\section{First-pass results}\label{sec:results}

\subsection{The comparative anatomy}\label{sec:results-anatomy}

Table~\ref{tab:ms-anatomy} is the paper's central object: the five
episodes side by side under the common lens. Three regularities
organise everything that follows. First, \emph{the negative shocks were regressive and the positive ones
were not progressive}: the gross energy shock would have cost the
bottom decile \MsEnergyGrossPctSpendDecOne\ per cent of its spending
against \MsEnergyGrossPctSpendDecTen\ per cent at the top (before the
policy response and demand adjustment), the TCA food shock \MsTcaWindowAvgPctSpendDecOne\ against
\MsTcaWindowAvgPctSpendDecTen\ per cent of spending on the
window-average path (\MsTcaPctSpendDecOne\ against
\MsTcaPctSpendDecTen\ at end-state), while the India agreement's consumer
gains are spread almost exactly proportionally
(\MsCetaBottomPctSpend\ versus \MsCetaTopPctSpend\ per cent of
spending) and are two orders of magnitude smaller. Second,
\emph{regressivity follows the affected category's budget-share
gradient}: food takes \MsTcaBudgetSharePctDecOne\ per cent of
bottom-decile spending against \MsTcaBudgetSharePctDecTen\ at the
top, energy \MsEnergyBudgetSharePctDecOne\ against
\MsEnergyBudgetSharePctDecTen, while clothing and footwear---where
the CETA gains land---are flat across the distribution: the
composition of a shock, not just its size, determines who bears it. Third, \emph{the cushioning came
from different parts of the state}: the one earnings-led episode was
cushioned automatically (\MsUSTariffDisplacementCushionPct--%
\MsUSTariffWageCutCushionPct\ per cent through the tax--benefit
system, from the companion study), while the price-led energy episode
was cushioned by discretion---the Energy Price Guarantee absorbed
\MsEnergyCushionSharePct\ per cent of the gross shock---and the
automatic price-side mechanism, benefit indexation, arrived a year
late (Section~\ref{sec:results-rulebook}).

\begin{table}[htbp]
\centering
\caption{Five shocks under one lens: first-pass comparability}
\label{tab:ms-anatomy}
\small
\begin{tabular}{llrrrl}
\toprule
Episode & Type, channel & Gross & \multicolumn{2}{c}{\% of spending} & Cushioning \\
\cmidrule(lr){4-5}
 & & shock & D1 & D10 & \\
\midrule
TCA food, 2021--23 & policy, prices & \pounds\MsTcaWindowAvgAggBn\,bn & \MsTcaWindowAvgPctSpendDecOne & \MsTcaWindowAvgPctSpendDecTen & in estimate \\
Energy, 2022--23 & market, prices & \pounds\MsEnergyGrossAggBn\,bn & \MsEnergyGrossPctSpendDecOne & \MsEnergyGrossPctSpendDecTen & EPG \MsEnergyCushionSharePct\% \\
US tariffs, 2025--26 & policy, earnings & \pounds\MsUSTariffGrossShockMillions\,m & \multicolumn{2}{c}{zero} & auto.\ \MsUSTariffDisplacementCushionPct--\MsUSTariffWageCutCushionPct\% \\
India CETA, 2026-- & policy, prices & \pounds\MsCetaTotalMillions\,m$^{g}$ & \MsCetaBottomPctSpend & \MsCetaTopPctSpend & n/a \\
CPTPP, 2024-- & benchmark & \pounds\MsCptppGdpBn\,bn$^{g}$ & \multicolumn{2}{c}{no structure} & n/a \\
\bottomrule
\end{tabular}
\begin{minipage}{.94\textwidth}\footnotesize
\emph{Notes:} $^{g}$~gain, not cost. Deciles are of equivalised
disposable income; entries are first-order incidence on published
expenditure-by-decile tables at fixed baskets. TCA on the
window-average path (end-state variant in the text); energy gross of
the policy response and demand adjustment. Episodes are typed
policy/market: the energy episode is trade-\emph{transmitted}, not a
policy act. ``Zero'' records that the UK did not retaliate, so no
consumer price vector exists. Cushioning entries measure different
mechanisms (a pass-through convention, a price subsidy, a tax--benefit
stabilisation coefficient) and are not summable; the common cushioning
estimand of Section~\ref{sec:framework} is computed at the
microsimulation stage, not here. Category budget-share gradients, which
generate these incidence patterns, are in
Figure~\ref{fig:budget-shares}; full per-decile detail in
Table~\ref{tab:decile-full}.
\end{minipage}
\end{table}

\begin{figure}[tbp]\centering
\includegraphics[width=\textwidth]{figures/fig_incidence_profile.png}
\caption{Household incidence by equivalised disposable income decile,
as a share of each decile's total spending. Costs are plotted as
positive; the India CETA series is a gain (dotted, open markers). The
energy episode is shown gross and net of the Energy Price Guarantee.}
\label{fig:incidence-profile}
\end{figure}

\begin{figure}[tbp]\centering
\includegraphics[width=\textwidth]{figures/fig_budget_shares.png}
\caption{Why each episode has the incidence it does: the affected
categories' shares of total spending by decile. Food and domestic
energy fall steeply with income; clothing and footwear---where the
India CETA gains land---are flat.}
\label{fig:budget-shares}
\end{figure}

% Generated by figures/make_figures.py -- do not edit by hand.
% Requires: \usepackage{booktabs}  (no other package needed)
% \rr = ragged-right inside a p{} cell, with \\ restored as the row
% terminator, so narrow columns neither overfull nor stretch-underfull.
\providecommand{\rr}{\raggedright\let\\\tabularnewline}

\begin{table}[tbp]
\centering
\small
\caption{Household incidence of the quantified episodes, by equivalised
disposable income decile. TCA food NTBs are on the window-average path
(0.30 of the end-state $+8\%$ effect); energy is the 2022--23 cap shock,
gross and net of the Energy Price Guarantee; the India CETA column is the
consumer gain at 100\% pass-through. Costs are positive.}
\label{tab:decile-full}
\begin{tabular}{@{}l rr rr rr r@{}}
\toprule
 & \multicolumn{2}{c}{TCA food NTBs} & \multicolumn{2}{c}{Energy, gross}
 & \multicolumn{2}{c}{Energy, net of EPG} & CETA gain \\
\cmidrule(lr){2-3}\cmidrule(lr){4-5}\cmidrule(lr){6-7}\cmidrule(l){8-8}
Decile & \pounds/yr & \% spend & \pounds/yr & \% spend & \pounds/yr & \% spend & \pounds/yr \\
\midrule
1 & 58 & 0.35 & 1,458 & 8.95 & 965 & 5.93 & 2.96 \\
2 & 63 & 0.35 & 1,504 & 8.35 & 996 & 5.53 & 3.93 \\
3 & 70 & 0.37 & 1,478 & 7.74 & 978 & 5.12 & 4.47 \\
4 & 74 & 0.33 & 1,631 & 7.21 & 1,080 & 4.78 & 3.90 \\
5 & 72 & 0.30 & 1,591 & 6.61 & 1,054 & 4.38 & 5.12 \\
6 & 84 & 0.30 & 1,584 & 5.62 & 1,049 & 3.72 & 5.37 \\
7 & 80 & 0.27 & 1,651 & 5.56 & 1,094 & 3.68 & 5.61 \\
8 & 84 & 0.26 & 1,631 & 4.99 & 1,080 & 3.31 & 6.56 \\
9 & 90 & 0.25 & 1,698 & 4.62 & 1,124 & 3.06 & 6.90 \\
10 & 94 & 0.20 & 1,986 & 4.17 & 1,315 & 2.76 & 9.93 \\
\midrule
D1/D10 ratio & --- & 1.80 & --- & 2.15 & --- & 2.15 & 0.98$^{\dagger}$ \\
\bottomrule
\end{tabular}

\vspace{2pt}
\begin{minipage}{\linewidth}\footnotesize
\emph{Notes.} \pounds\ per household per year; \% spend is the burden or
gain as a percentage of that decile's total expenditure. The EPG cushions
33.8\% of the gross energy shock at every decile by construction.
$^{\dagger}$~CETA ratio is on the \%-of-spending basis (\pounds\ gains are
not comparable across deciles). Spending bases: ONS Family Spending
workbook~1, sheet 3.1E, FYE2022 vintage (TCA, energy) and FYE2025 vintage (CETA).
\end{minipage}
\end{table}



\subsection{Episode detail}

\paragraph{TCA food prices.} The headline incidence uses the
window-average path implied by the published cumulative: the
\pounds\MsTcaPublishedCumulGbp-per-household estimate implies an
average effect of about \MsTcaImpliedAvgShare\ of the 8 per cent
end-state over the window, giving \pounds\MsTcaWindowAvgGbpDecOne\
(bottom decile) to \pounds\MsTcaWindowAvgGbpDecTen\ (top decile) a
year, \pounds\MsTcaWindowAvgAggBn\ billion a year in aggregate---
regressive as a share of spending
(\MsTcaWindowAvgPctSpendDecOne\ versus
\MsTcaWindowAvgPctSpendDecTen\ per cent) because food takes
\MsTcaBudgetSharePctDecOne\ per cent of bottom-decile spending
against \MsTcaBudgetSharePctDecTen\ at the top. The end-state
figures (\pounds\MsTcaAnnGbpDecOne--\MsTcaAnnGbpDecTen\ a year;
\pounds\MsTcaAggBnPerYear\ billion) are reported as a labelled
variant: the pipeline's linear-ramp cross-check against the published
cumulative comes in at ratio \MsTcaCrossCheckRatio, showing the price
effect was back-loaded, so end-state overstates the lived burden. The
microsimulation stage will carry the published path.

\paragraph{The energy shock and the EPG.} At cap levels, the gross
price shock (\pounds\MsEnergyCapBase\ to \pounds\MsEnergyCapPeak\
typical bill, +\MsEnergyGrossFactorPct\ per cent) \emph{would have}
cost the bottom decile \pounds\MsEnergyGrossGbpDecOne\ a year at
fixed baskets---\MsEnergyGrossPctSpendDecOne\ per cent of its total
spending, more than double the top decile's share. Three labelled
qualifications bound that number. A demand-response variant (unit
elasticity applied at half weight, consistent with the double-digit
weather-adjusted demand reductions observed in 2022--23) scales the
burden to \MsEnergyDemandScalePct\ per cent of the fixed-basket
figure (\pounds\MsEnergyGrossGbpDecOneDemand\ for the bottom decile;
\pounds\MsEnergyGrossAggBnDemand\ billion aggregate). An
imported-component bound, applying the \citet{dhingrapage2023}
finding that over \MsEnergyImportedShareBound\ per cent of the CPI
rise was imported, caps the trade-transmitted share of the gross
aggregate at roughly \pounds\MsEnergyGrossAggBnImported\ billion;
the remainder is domestic cost and margin. And the annualised
run-rate holds the October 2022 cap for a full year, which no
household experienced; the six-month window figure is
\pounds\MsEnergyCushionAggSixMonthsBn\ billion of EPG absorption
against the OBR's \pounds\MsEnergyObrEpgBn\ billion costing---a
context anchor, not a validation. Within those bounds the shape is
robust: the Energy Price Guarantee truncated the rise to
+\MsEnergyNetFactorPct\ per cent, absorbing
\pounds\MsEnergyCushionGbpDecOne\ a year for a bottom-decile
household and \MsEnergyCushionSharePct\ per cent of the gross shock.
Because the EPG capped the \emph{price}, its cushioning share is
constant across deciles by construction. Two qualifications follow
from the frontier evidence. First, decile means understate the
problem the instrument faced: \citet{levell2025} show that income
explains only a small share of the variation in household energy
spending, so exposure within a decile is far more dispersed than the
gradient between deciles---an argument for the household-record stage,
where that dispersion is observable, and against reading any
decile-mean cushioning share as a statement about individual
households. Second, a proportional instrument is not thereby a
mistaken one: \citet{levell2025} find that an optimal package retains
a strictly positive price subsidy even when transfers can condition on
income and past usage, precisely because income targets energy
exposure so poorly. The finding here is descriptive---the cap
cushioned proportionally---not a welfare verdict. Whether the overall
discretionary response was progressive cannot be judged from the EPG
alone: the Energy Bills Support Scheme's flat \pounds 400, the
means-tested cost-of-living payments and the council-tax rebate---
jointly of the same order as the EPG and deliberately targeted---are
not modelled in this first pass, and the distributional verdict on
the discretionary stack is deferred to the microsimulation stage.

\begin{figure}[tbp]\centering
\includegraphics[width=\textwidth]{figures/fig_energy_epg.png}
\caption{The 2022--23 energy shock by decile: the gross cap-level
burden split into the part borne by households net of the Energy Price
Guarantee and the part the Guarantee absorbed. Because the instrument
capped a price, its cushioning share is constant across the
distribution by construction.}
\label{fig:energy-epg}
\end{figure}

\paragraph{The US tariffs.} The UK consumer side is zero---no
retaliation occurred---so this episode's incidence is entirely
earnings-led: a gross loss of \pounds\MsUSTariffGrossShockMillions\
million a year cushioned by the automatic tax--benefit system at
\MsUSTariffDisplacementCushionPct\ per cent (displacement) to
\MsUSTariffWageCutCushionPct\ per cent (wage cuts), the companion
study's headline carried over as this episode's second stage.

\paragraph{The India gains, and the placebo.} At full pass-through the
CETA duty cuts are worth \pounds\MsCetaGainGbpPtHundredDecOne\ a year
to a bottom-decile household and \pounds\MsCetaGainGbpPtHundredDecTen\
to a top-decile one---\MsCetaVerdict\ as a share of spending, though
pro-rich in pounds (the top decile gains about three times as much in
cash), and economically negligible at any pass-through dial
(\pounds\MsCetaGainGbpPtFiftyDecOne--\MsCetaGainGbpPtFiftyDecTen\ at
50 per cent). The CPTPP row is a near-zero benchmark, included as a
data-limitation case rather than a placebo test: its
\pounds\MsCptppGdpBn\ billion long-run GDP estimate divides to
\pounds\MsCptppNaiveGbpPerHousehold\ per household per year with no
defensible distributional structure, and the pipeline reports exactly
that and nothing more.

\subsection{The rulebook seed results}\label{sec:results-rulebook}

The vintage axis is seeded with two statutory facts. First, the
uprating lag: Universal Credit's standard allowance rose 3.1 per cent
in April 2022 (indexed to the previous September's inflation) while
prices rose at double digits. Against month-by-month contemporaneous
indexation---an upper-bound counterfactual no benefit system
operates---a single adult's standard allowance fell
\pounds\MsUcLagCostGbp\ short over 2022--23
(\MsUcLagCostPctAllowance\ per cent of the annual allowance); against
the policy-feasible alternative of a single April 2022 uprating at
contemporaneous annual CPI, \pounds\MsUcLagCostFlatAltGbp. Both
figures are for the single-adult standard allowance only, before
elements and caseload weighting; the range, not either point, is the
seed result. Neither the lag nor its cost is a discovery of this
paper: \citet{cribb2023} quantified the 2022--23 shortfall and scored
the discretionary cost-of-living payments against a timely-indexation
counterfactual, and \citet{obrwelfare2022} put the aggregate real fall
at roughly \pounds 12 billion---a figure the household-record stage
should reproduce as a validation target. What is new here is the
join: the lag priced against trade-shock incidence, on the same
households, across episodes. The lag also does not wash out over the
cycle: \citet{judgemurphy2023} count ten of fifteen years in which
working-age benefits rose by less than outturn inflation. The shortfall
is lost to timing alone, at exactly the moment the terms-of-trade shock
peaked. Figure~\ref{fig:uprating-lag} shows the gap month by month. Second,
sequencing: the \pounds\MsUcUpliftRemovalGbp-a-year
Covid uplift was removed in October 2021, months before the energy
shock landed, so the lowest-income households entered the largest
price-led trade shock of the period from a rulebook trough. Together
with the EPG's absorption of \MsEnergyCushionSharePct\ per cent of the
gross price shock, the seed picture of the British response is:
\emph{automatic on the earnings margin, discretionary on the price
margin, and lagged precisely where the price shocks bit hardest}---with
the progressivity of the full discretionary stack an open question for
the microsimulation stage. The microsimulation
stage turns these representative-household facts into
population-weighted decompositions.

\begin{figure}[tbp]\centering
\includegraphics[width=\textwidth]{figures/fig_uprating_lag.png}
\caption{The uprating lag in the crisis year. The Universal Credit
standard allowance for a single adult aged 25 or over (statutory path,
flat from April 2022) against a counterfactual indexed to
contemporaneous CPI. The shaded area is the shortfall, summing to
\pounds\MsUcLagCostGbp\ over 2022--23.}
\label{fig:uprating-lag}
\end{figure}
